\documentclass[aspectratio=169,11pt]{beamer}

% Portable on purpose: no shell-escape, custom fonts, or external theme packages.
\usetheme{Madrid}
\usecolortheme{default}
\setbeamertemplate{navigation symbols}{}

\usepackage[T1]{fontenc}
\usepackage[utf8]{inputenc}
\usepackage{listings}
\usepackage{xcolor}
\usepackage{graphicx}
\usepackage{booktabs}

% Keep code readable from the back of the room.
\lstdefinestyle{python}{
  language=Python,
  basicstyle=\ttfamily\small,
  showstringspaces=false,
  columns=fullflexible,
  keepspaces=true,
  breaklines=true,
  frame=single,
  tabsize=4
}
\lstset{style=python}

\title{Presentation Title}
\subtitle{Optional subtitle}
\author{Jeremy Evert}
\institute{SWOSU}
\date{}

\begin{document}

\begin{frame}
  \titlepage
\end{frame}

% Question slide: one idea, lots of air.
\begin{frame}{A useful question}
  \centering
  \Huge What is this thing, and what can it do?
\end{frame}

% Code slide: add [fragile] whenever listings/verbatim appear.
\begin{frame}[fragile]{Tiny code, big idea}
\begin{lstlisting}
text = "hello"
print(text.upper())
\end{lstlisting}
\end{frame}

% Comparison / concept slide.
\begin{frame}{Three pieces}
  \begin{columns}[T,onlytextwidth]
    \column{0.31\textwidth}
      \centering
      \Large Type

    \column{0.31\textwidth}
      \centering
      \Large State

    \column{0.31\textwidth}
      \centering
      \Large Behavior
  \end{columns}
\end{frame}

\begin{frame}{Exit question}
  \centering
  \LARGE What object are you holding, and what can it do?
\end{frame}

\end{document}
